% SGA 5，第 I 讲附录，原 PDF 第 67--84 页。
\SourcePage{67}
\SGAsection{附录}{附录}
\begin{center}L. Illusie 撰\end{center}
前面的理论很容易推广到比常值环 $\mathbf Z/n\mathbf Z$ 更一般的系数环。我们将简要说明：当系数环是\Term{locally-constant}{局部常值}的、被某个 $n>0$ \Term{annihilate}{零化}，而且其\Term{geometric-fiber}{几何纤维}是\Term{noetherian-ring}{诺特环}时，理论应当如何表述。不过，这还只是一个暂时的框架\footnote{据 Grothendieck 所说！}：迟早还须将它整体化，即在\Term{prescheme}{预概形} $X$ 上，以\Term{ringed-site}{环化景} $(X_{\text{ét}},\mathbf Z/n\mathbf Z)$ 上的相对\Term{prescheme}{预概形}作为“系数”。

\SGAsection{1}{可构造层}
\textbf{1.1.} 设 $X$ 为\Term{prescheme}{预概形}，$A_X$ 为 $X$ 上（对于\Term{etale}{平展}\Term{topology}{拓扑}）的\Term{locally-constant}{局部常值}\Term{ringed-sheaf}{环层}，其\Term{geometric-fiber}{几何纤维}为\Term{left-noetherian-ring}{左诺特环}。以 $\mathsf A^X$ 表示左 $A_X$ \Term{module-category}{模范畴}，以 $D(\mathsf A^X)$（在不致混淆时也记作 $D(X)$）表示相应的\Term{derived-category}{导出范畴}。\Term{constructible}{可构造}（左）$A_X$ 模的定义与（SGA 4 IX 2.3）相同，只须把其中“\Term{locally-constant}{局部常值}且其值为\Term{finitely-presented}{有限呈示}的”替换为“作为 $A_{U_i}$ 模是\Term{finitely-presented}{有限呈示}的”。
以 $D_c(X)$ 表示 $D(X)$ 的如下\Term{full-subcategory}{全子范畴}：其\Term{object}{对象}为对每个 $i$ 都使 $H^i(E)$ 为\Term{constructible}{可构造} $A_X$ 模的\Term{complex}{复形} $E$。它是 $D(X)$ 的\Term{triangulated-category}{三角子范畴}。

\SourcePage{68}
\textbf{1.2.} 设 $E\in\operatorname{ob}D(X)$。若对每个 $i$，$H^i(E)$ 都是\Term{finitely-presented}{有限呈示}的 $A_X$ 模（也可以说是\Term{locally-constant}{局部常值}\Term{constructible}{可构造}模），则称 $E$ 为\Term{pseudo-coherent}{伪凝聚}的（参见第 II 讲 3.11）。以 $D(X)_{\mathrm{coh}}$（或 $D_{\mathrm{coh}}(X)$）表示由\Term{pseudo-coherent}{伪凝聚}\Term{object}{对象}组成的 $D(X)$ 的\Term{full-subcategory}{全子范畴}。由（SGA 4 IX 2.1），它是 $D_c(X)$ 的\Term{triangulated-category}{三角子范畴}。

\begin{SGAresult}{命题}{1.3（\Term{resolution}{分解}\Term{lemma}{引理}）}
设 $X$ 为\Term{quasi-compact}{拟紧}且\Term{quasi-separated}{拟分离}的（分别地，\Term{noetherian}{诺特}的）\Term{prescheme}{预概形}，$C$ 为\Term{constructible}{可构造}左 $A_X$ \Term{module-category}{模范畴} $\operatorname{Cons}(X)$ 的一个\Term{strictly-full-subcategory}{严格全子范畴}。假定 $C$ 对\Term{group-extension}{扩张}\Term{stable}{稳定}。下列条件\Term{equivalence}{等价}：
\begin{enumerate}
\item[(i)] $C=\operatorname{Cons}(X)$；
\item[(ii)] $C$ 包含所有形如 $i_!(F)$ 的层，其中 $i:Y\to X$ 为\Term{finitely-presented}{有限呈示}的\Term{subprescheme}{子预概形}（分别地，\Term{integral}{整}的\Term{subprescheme}{子预概形}），$F$ 在 $Y$ 上是\Term{locally-constant}{局部常值}\Term{constructible}{可构造}的。
\end{enumerate}
若 $A_X$ 为\Term{constant-sheaf}{常值环层}，其值为被某个整数 $n>0$ \Term{annihilate}{零化}的\Term{left-noetherian-ring}{左诺特环} $A$，则上述条件又\Term{equivalence}{等价}于：
\begin{enumerate}
\item[(iii)] $C$ 对\Term{direct-summand}{直和因子}\Term{stable}{稳定}，并包含所有形如 $f_*i_!(M_U)$ 的层，其中 $i:U\to Y$ 为\Term{open-immersion}{开浸入}（分别地，为某个\Term{integral}{整}\Term{prescheme}{预概形}的\Term{open-subset}{开集}），$f:Y\to X$ 为\Term{finite-morphism}{有限态射}（分别地，$f$ \Term{finite}{有限}，且 $Y$ 的\Term{generic-point}{一般点}在其像上\Term{separable}{可分}），而 $M_U$ 为常值左 $A_U$ 模，其值为被 $n$ 的某个素因子\Term{annihilate}{零化}的\Term{cyclic-module}{单生成} $A$ 模。
\end{enumerate}
最后，若 $X$ 为\Term{noetherian}{诺特}的且满足条件 $\mathrm{(reg)}$（第 I 讲 3.1.1），则条件 (ii)（分别地，(iii)）\Term{equivalence}{等价}于进一步要求 $Y$（分别地，$U$）\Term{regular}{正则}时的对应条件。
\end{SGAresult}
\SourcePage{69}
\begin{proof}
条件 (i) 与 (ii) 的\Term{equivalence}{等价}性，证明同 $A_X$ 为常值时一样，参见（SGA 4 IX 2.4、2.5）。

假定 $A_X$ 是值为 $A$ 的\Term{constant-sheaf}{常值环层}，且 $A$ 被 $n$ \Term{annihilate}{零化}，来证明 (iii) 蕴涵 (ii)。取 $E\in\operatorname{ob}\operatorname{Cons}(X)$，要证 $E\in\operatorname{ob}C$。由 (i)$\Longleftrightarrow$(ii)，可以假定 $E=j_!(F)$，其中 $j:Z\to X$ 为\Term{finitely-presented}{有限呈示}（分别地，\Term{integral}{整}）\Term{subprescheme}{子预概形}的\Term{immersion}{浸入}，而 $F$ 为\Term{locally-constant}{局部常值}\Term{constructible}{可构造} $A_Z$ 模；再作少许\Term{resolution}{分解}，还可假定 $F$ 由一个群为 $G$ 的\Term{etale-covering}{平展覆盖} $g:Z'\to Z$ \Term{trivialization}{平凡化}。因此，层 $F$ 由一个带 $G$ \Term{action}{作用}的\Term{finite-type}{有限型} $A$ 模 $M$ 定义。若 $\ell$ 为 $n$ 的素因子，则 $M$ 的 $\ell$ \Term{torsion-subgroup}{挠子群}是 $M$ 的一个 $G$-$A$ \Term{submodule}{子模} $M_\ell$，且有
\[
M=\bigoplus_\ell M_\ell,
\]
其中 $\ell$ 遍历 $n$ 的所有素因子。由于 $C$ 对\Term{direct-summand}{直和因子}\Term{stable}{稳定}，只须处理 $M$ 为 $\ell$ \Term{torsion-module}{挠模}的情形。
取 $G$ 的一个 Sylow $\ell$ \Term{subgroup}{子群} $H$，以 $h:U=Z'/H\to Z$ 表示对应于 $H$ 的中间\Term{covering}{覆盖}。由于 $h$ 的\Term{degree}{次数}与 $\ell$ 互素，“迹方法”（SGA 4 IX 5.1）表明\Term{adjoint}{伴随箭头}
\[
F\longrightarrow h_*h^*(F)
\]
将 $F$ \Term{isomorphism}{同构}到一个\Term{direct-summand}{直和因子}上，所以只须证明 $j_!h_*h^*(F)\in\operatorname{ob}C$。由“主\Term{theorem}{定理}”（EGA IV 8.12.6），存在\Term{commutative-diagram}{交换图}
\[
\begin{tikzcd}
U\arrow[r,"i"]\arrow[d,"h"']&Y\arrow[d,"f"]\\
Z\arrow[r,"j"']&X,
\end{tikzcd}
\]
\SourcePage{70}
其中 $i$ 为\Term{open-immersion}{开浸入}，$f$ 为\Term{finite-morphism}{有限态射}。有
\[
j_!h_*h^*(F)\simeq f_*i_!h^*(F).
\]
而 $h^*(F)$ 由视作 $H$-$A$ 模的 $M$ 定义。因此，蕴涵 (iii)$\Longrightarrow$(ii) 将由下面的\Term{lemma}{引理}得到。
\end{proof}

\begin{SGAresult}{\Term{lemma}{引理}}{1.3.1}
设 $A$ 为\Term{left-noetherian-ring}{左诺特环}，$\ell$ 为素数，$H$ 为\Term{finite}{有限} $\ell$ 群，$M$ 为被 $\ell^\nu$（$\nu\geq1$）\Term{annihilate}{零化}的\Term{finite-type}{有限型}左 $H$-$A$ 模。则 $M$ 存在由 $H$-$A$ \Term{submodule}{子模}组成的\Term{finite}{有限}\Term{filtered}{滤过}，使各相继商均为被 $\ell$ \Term{annihilate}{零化}且 $H$ 在其上\Term{trivial-action}{平凡作用}的\Term{cyclic-module}{单生成} $A$ 模。
\end{SGAresult}
\begin{proof}
$H$-$A$ 模的\Term{exact-sequence}{正合列}
\[
0\longrightarrow\ell M\longrightarrow M\longrightarrow M/\ell M\longrightarrow0
\]
使我们可对 $\nu$ 归纳，归结到 $M$ 被 $\ell$ \Term{annihilate}{零化}的情形。
现证 $M$ 存在由 $H$-$A$ \Term{submodule}{子模}组成的\Term{finite}{有限}\Term{filtered}{滤过}，使 $H$ 在各相继商上\Term{trivial-action}{平凡作用}。在 $M$ 中具有上述性质的 $H$-$A$ \Term{submodule}{子模}所成的\Term{set}{集合}上作\Term{noetherian-induction}{诺特归纳}，就归结为证明：$M\ne0$ 蕴涵 $M^H\ne0$。容易验证，\Term{augmentation}{增广} $\mathbf F_\ell[H]\to\mathbf F_\ell$ 的核是一个\Term{nilpotent-ideal}{幂零理想}，故 \Term{nakayama-lemma}{Nakayama 引理}说明，每个非零的 $H$-$\mathbf F_\ell$ 模都有非零\Term{invariant}{不变}向量，这就证明了所需断言。
因此，$M$ 存在\Term{finite}{有限}\Term{filtered}{滤过}，其相继商是被 $\ell$ \Term{annihilate}{零化}且 $H$ 在其上\Term{trivial-action}{平凡作用}的\Term{finite-type}{有限型} $A$ 模。再作\Term{noetherian-induction}{诺特归纳}，甚至可以使相继商为\Term{cyclic-module}{单生成} $A$ 模，从而完成（1.3.1）的证明。
\end{proof}
（1.3）的最后一项断言是直接的，留给读者证明。

\SourcePage{71}
\SGAsection{2}{拟内射维数}
沿用（1.1）的记号，仿照（第 I 讲 1.1、1.4）定义\Term{quasi-injective}{拟内射} $A_X$ 模以及 $D(X)$ 的\Term{object}{对象}的\Term{quasi-injective-dimension}{拟内射维数}。命题（第 I 讲 1.3、1.5、1.6）原样推广。

有时引入逐点的\Term{quasi-injective-dimension}{拟内射维数}概念很方便。设 $x\in X$，$K\in\operatorname{ob}D^+(X)$。若对每个\Term{constructible}{可构造} $A_X$ 模 $F$ 均有
\[
\mathcal Ext^i(F,K)_{\bar x}=0\qquad(i>N),
\]
则称 $K$ 在 $\bar x$ 处的\Term{quasi-injective-dimension}{拟内射维数} $\leq N$。将满足上述条件的 $N$ 的\Term{infimum}{下确界}称为 $K$ 在 $\bar x$ 处的\Term{quasi-injective-dimension}{拟内射维数}，记作 $\operatorname{dim.q.inj}_{\bar x}(K)$。当然有
\[
\operatorname{dim.q.inj}(K)=\sup_{\bar x}\operatorname{dim.q.inj}_{\bar x}(K),
\]
其中 $\bar x$ 遍历 $X$ 的所有\Term{geometric-point}{几何点}。

把上述概念与 Verdier 在 [1] 中引入的\Term{strict-topological-dimension}{严格拓扑维数}概念作比较很有意义。若对 $X$ 的每个\Term{subprescheme}{子预概形} $Y$ 均有
\[
\mathcal Ext^i(A_Y,K)_{\bar x}=\mathcal H_Y^i(K)_{\bar x}=0\qquad(i>N),
\]
则称 $K\in\operatorname{ob}D^+(X)$ 在 $X$ 的\Term{geometric-point}{几何点} $\bar x$ 处的\Term{strict-topological-dimension}{严格拓扑维数} $\leq N$。将满足上述条件的 $N$ 的\Term{infimum}{下确界}称为 $K$ 在 $\bar x$ 处的\Term{strict-topological-dimension}{严格拓扑维数}，记作 $\operatorname{dimstop}_{\bar x}(K)$。

\SourcePage{72}
\Term{quasi-injective-dimension}{拟内射维数}先验地依赖于系数环，而\Term{strict-topological-dimension}{严格拓扑维数}只依赖于底层\Term{abelian-sheaf}{阿贝尔层}\Term{complex}{复形}。显然有
\[
\operatorname{dimstop}_{\bar x}(K)\leq\operatorname{dim.q.inj}_{\bar x}(K).
\]
下文将给出等号成立的例子。

\SGAsection{3}{对偶化复形：定义与形式性质}
保持前面的记号，但假定 $A_X$ 交换。\Term{bidualizing-pair}{双对偶化二元组}和\Term{dualizing-complex}{对偶化复形}的概念仿照（第 I 讲 1.7）定义。断言（第 I 讲 1.9、1.11、1.14）无须修改仍然成立。断言（第 I 讲 1.12、1.13、1.15）也成立，但须假定 $A$ 被某个与 $Y$ 的\Term{residue-characteristic}{剩余特征}互素的整数 $n>0$ \Term{annihilate}{零化}。最后，只要假定对 $X$ 的每个\Term{geometric-point}{几何点} $\bar x$，$\operatorname{Spec}(A_{\bar x})$ 都\Term{connected}{连通}，则唯一性\Term{theorem}{定理}（第 I 讲 2.1）也成立。

\SGAsection{4}{绝对纯性}
回到\Term{absolute-cohomological-purity}{绝对上同调纯性}的情形（第 I 讲 3.1.4）。设 $X$ 为\Term{locally-noetherian}{局部诺特}\Term{regular}{正则}\Term{prescheme}{预概形}，$n>0$ 为与 $X$ 的\Term{residue-characteristic}{剩余特征}互素的整数。设 $i:Y\to X$ 为\Term{regular}{正则}\Term{subprescheme}{子预概形}，在每一点处的\Term{codimension}{余维}均为 $d>0$。回顾一下，称二元组 $(Y,X)$ 满足\Term{absolute-cohomological-purity}{绝对上同调纯性}，是指“\Term{local-fundamental-class}{局部基本类}”\Term{homomorphism}{同态}（SGA $4\frac12$ Cycle 2.2）
\SourcePage{73}
\[
(\mu_n)_Y^{\otimes-d}[-2d]\longrightarrow Ri^!(\mathbf Z/n\mathbf Z)_X
\tag{4.1}
\]
为\Term{isomorphism}{同构}。

现在给定 $X$ 上被 $n$ \Term{annihilate}{零化}的\Term{ringed-sheaf}{环层} $A_X$。暂记 $G=(\mathbf Z/n\mathbf Z)_X$。则对每个 $F\in\operatorname{ob}D^+(\mathsf A^X)$，可定义 $D^+(\mathsf A^X)$ 中的\Term{canonical}{典范}\Term{arrow}{箭头}
\[
i^*(F)\mathbin{\mathop\otimes\limits^L}_{\mathbf Z/n\mathbf Z}Ri^!(G)
\longrightarrow Ri^!\bigl(F\mathbin{\mathop\otimes\limits^L}_{\mathbf Z/n\mathbf Z}G\bigr).
\tag{4.2}
\]
事实上，由\Term{duality}{对偶性}（SGA 4 XVIII 3.1.6），这\Term{equivalence}{等价}于定义一个\Term{arrow}{箭头}
\[
i_!\bigl(i^*(F)\mathbin{\mathop\otimes\limits^L}_{\mathbf Z/n\mathbf Z}Ri^!(G)\bigr)
\longrightarrow F\mathbin{\mathop\otimes\limits^L}_{\mathbf Z/n\mathbf Z}G;
\]
我们取 Künneth \Term{arrow}{箭头}（SGA 4 XVII 5.4.2.2）
\[
i_!\bigl(i^*(F)\mathbin{\mathop\otimes\limits^L}_{\mathbf Z/n\mathbf Z}Ri^!(G)\bigr)
\longrightarrow F\mathbin{\mathop\otimes\limits^L}_{\mathbf Z/n\mathbf Z}i_!Ri^!(G)
\]
与由\Term{trace-homomorphism}{迹同态}（SGA 4 XVIII 3.1.6）\Term{derive}{导出}的\Term{arrow}{箭头}
\[
F\mathbin{\mathop\otimes\limits^L}_{\mathbf Z/n\mathbf Z}i_!Ri^!(G)
\longrightarrow F\mathbin{\mathop\otimes\limits^L}_{\mathbf Z/n\mathbf Z}G
\]
之复合。
应当注意，定义（4.2）时，我们只使用了 $G$ 和 $Ri^!(G)$ 在 $\mathbf Z/n\mathbf Z$ 上具有\Term{finite}{有限}\Term{tor-dimension}{Tor 维数}这一事实。另一方面，假定\Term{functor}{函子} $i^!$ 在 $(\mathbf Z/n\mathbf Z)_X$ \Term{module-category}{模范畴}上具有\Term{finite}{有限}\Term{cohomological-dimension}{上同调维数}：例如，当 $X$ 在域上为\Term{finite-type}{有限型}时（SGA 4 XIV 3.1），或更一般地，当 $X$ 为\Term{equal-characteristic}{等特征}\Term{excellent}{优良}\Term{prescheme}{预概形}且可用\Term{resolution-of-singularities}{奇点消解}时（SGA 4 XIX 6.1），或当 $X$ 在 $\operatorname{Spec}(\mathbf Z)$ 上为\Term{finite-type}{有限型}时（SGA 4 X 6），均属这种情形。此时，对每个 $F\in\operatorname{ob}D(\mathsf A^X)$ 都可以定义\Term{arrow}{箭头}（4.2）。

因此，复合（4.1）与（4.2），就得到对 $F\in\operatorname{ob}D^+(\mathsf A^X)$ 的\Term{canonical}{典范}\Term{arrow}{箭头}
\SourcePage{74}
\[
i^*(F)\mathbin{\mathop\otimes\limits^L}_{\mathbf Z/n\mathbf Z}(\mu_n)_Y^{\otimes-d}[-2d]
\longrightarrow Ri^!(F).
\tag{4.3}
\]
假定 $A_X$ 为\Term{locally-constant}{局部常值}的，且 $F\in\operatorname{ob}D^+(\mathsf A^X)$ 满足：对每个 $i$，$H^i(F)$ 均为\Term{locally-constant}{局部常值} $A_X$ 模。则\textbf{若 $i^!$ 具有\Term{finite}{有限}\Term{cohomological-dimension}{上同调维数}，\Term{arrow}{箭头}（4.2）就是\Term{isomorphism}{同构}；因而若 $(Y,X)$ 满足\Term{purity}{纯性}，\Term{arrow}{箭头}（4.3）也是\Term{isomorphism}{同构}。}

事实上，借助\Term{resolution}{分解}，首先归结到 $F$ 集中于\Term{degree}{次数} $0$ 的情形。由于问题是\Term{local-scheme}{局部}的，可以假定 $X$ 为\Term{affine}{仿射}的，且 $A_X$、$F$ 都是常值的。用自由 $A_X$ 模向\Term{left-resolution}{左消解} $F$，再利用 $i^!$ 的\Term{cohomological-dimension}{上同调维数}\Term{finite}{有限}这一假设，就归结到 $F$ 为\Term{free-module}{自由模}的情形。由直接的取\Term{limit}{极限}过程，又归结到 $F$ 为\Term{finite-type}{有限型}\Term{free-module}{自由模}，最终归结到 $F=A_X$。所以，只须在 $A_X=(\mathbf Z/n\mathbf Z)_X$ 时证明断言。但前面的\Term{resolution}{分解}将此归结到 $F=(\mathbf Z/n\mathbf Z)_X$，而此时（4.2）就是恒等\Term{arrow}{箭头}。

\begin{SGAresult}{定义}{4.4}
给定\Term{ringed-sheaf}{环层} $A_X$（\Term{locally-constant}{局部常值}、被 $n$ \Term{annihilate}{零化}），若对每个使 $H^i(F)$ 均为\Term{locally-constant}{局部常值}的 $F\in\operatorname{ob}D^+(\mathsf A^X)$，\Term{canonical}{典范}\Term{arrow}{箭头}（4.3）都是\Term{isomorphism}{同构}，则称二元组 $(Y,X)$ 满足\Term{strong-cohomological-purity}{强意义的上同调纯性}。

若\Term{local-scheme}{局部}地（对于\Term{etale}{平展}\Term{topology}{拓扑}），每个如上的二元组 $(Y,X)$ 都满足\Term{strong-cohomological-purity}{强意义的纯性}，则称 $X$ 上的\Term{purity-theorem}{纯性定理}在强意义下成立。
\end{SGAresult}
由前面的讨论，要使 $(Y,X)$ 满足\Term{strong-cohomological-purity}{强意义的纯性}，只须 $(Y,X)$ 满足通常意义的\Term{purity}{纯性}（第 I 讲 3.1），且\Term{functor}{函子} $i^!$ 具有\Term{finite}{有限}\Term{cohomological-dimension}{上同调维数}。

\SourcePage{75}
\SGAsection{5}{拟内射维数的上界}
\begin{SGAresult}{命题}{5.1}
设 $X$ 为\Term{locally-noetherian}{局部诺特}\Term{regular}{正则}\Term{prescheme}{预概形}，配备一个\Term{locally-constant}{局部常值}\Term{ringed-sheaf}{环层} $A_X$，它被某个与 $X$ 的\Term{residue-characteristic}{剩余特征}互素的整数 $n>0$ \Term{annihilate}{零化}，且其\Term{geometric-fiber}{几何纤维}为\Term{left-noetherian-ring}{左诺特环}。假定 $X$ 满足条件 $\mathrm{(reg)}$（第 I 讲 3.1.1）和 $\mathrm{(cdloc)}_n$（第 I 讲 3.1.3），并且 $X$ 上的\Term{purity-theorem}{纯性定理}在强意义下成立（4.4）。（当 $X$ 在某个完全域上\Term{smooth}{光滑}，或 $X$ 为零\Term{characteristic}{特征}\Term{excellent}{优良}\Term{prescheme}{预概形}时，这些条件都满足。）设 $F\in\operatorname{ob}D^+(\mathsf A^X)$，其 $H^i(F)$ 均为\Term{locally-constant}{局部常值}的；设 $x$ 为 $X$ 的一点，满足 $F_{\bar x}\ne0$ 且 $\operatorname{dim.inj}_{A_{\bar x}}(F_{\bar x})<+\infty$。则有
\[
\operatorname{dim.q.inj}_{\bar x}(F)
=2\dim(\mathcal O_{X,x})+\operatorname{dim.inj}(F_{\bar x}),
\]
其中使用第 2 节的记号。
\end{SGAresult}
\begin{proof}
只须逐步重作（第 I 讲 3.2.1）的证明。先证不等式
\[
\operatorname{dim.q.inj}_{\bar x}(F)
\leq2\dim(\mathcal O_{X,x})+\operatorname{dim.inj}(F_{\bar x})=N,
\]
也就是关系
\[
\mathcal Ext^i(E,F)_{\bar x}=0
\quad\text{当 }i>N\text{，且 }E\text{ 为任意可构造 }A_X\text{ 模}.
\tag{*}
\]
由于问题是\Term{local-scheme}{局部}的，借助\Term{resolution}{分解}（1.3），只须对 $E=i_!(M)$ 验证（*），其中 $i:Y\to X$ 为\Term{integral}{整}\Term{regular}{正则}\Term{subprescheme}{子预概形}，它在某个\Term{open-subset}{开集} $D(f)$（$f\in\Gamma(X,\mathcal O_X)$）中为闭的，在 $D(f)$ 的每一点处\Term{codimension}{余维}均为 $d$，$M$ 为\Term{locally-constant}{局部常值}\Term{constructible}{可构造} $A_Y$ 模（$A_Y=i^*A_X$）。
\SourcePage{76}
有\Term{duality-isomorphism}{对偶同构}（SGA 4 XVIII\textup{〔原文未列具体条目〕}）
\[
R\mathcal Hom(i_!(M),F)\simeq Ri_*R\mathcal Hom(M,Ri^!F).
\]
另一方面，由\Term{purity-theorem}{纯性定理}有\Term{isomorphism}{同构}（4.3）
\[
i^*(F)\mathbin{\mathop\otimes\limits^L}_{\mathbf Z/n\mathbf Z}(\mu_n)_Y^{\otimes-d}[-2d]
\xrightarrow{\sim}Ri^!(F).
\]
设 $\bar y$ 为 $Y$ 的\Term{geometric-point}{几何点}。由于 $M$ 为\Term{locally-constant}{局部常值}\Term{constructible}{可构造}的，\Term{canonical}{典范}地有
\[
R\mathcal Hom(M,Ri^!F)_{\bar y}
\simeq R\operatorname{Hom}\bigl(M_{\bar y},(Ri^!F)_{\bar y}\bigr),
\]
故由（4.3）得
\[
R\mathcal Hom(M,Ri^!F)_{\bar y}
\simeq R\operatorname{Hom}(M_{\bar y},F_{\bar y})[-2d].
\]
另一方面，可以假定我们已处在 $x$ 的一个\Term{connected}{连通}\Term{neighborhood}{邻域}内。于是有
\[
F_{\bar x}\simeq F_{\bar y},\qquad A_{\bar x}\simeq A_{\bar y},
\]
因为 $A_X$ 和各个 $H^i(F)$ 都是\Term{locally-constant}{局部常值}的；由此得
\[
\operatorname{Ext}^i(M_{\bar y},F_{\bar y})=0
\qquad\bigl(i>\operatorname{dim.inj}(F_{\bar x})\bigr),
\]
最终得到
\[
\mathcal Ext^q(M,Ri^!F)=0
\qquad\bigl(q>2d+\operatorname{dim.inj}(F_{\bar x})\bigr).
\]
置 $G=R\mathcal Hom(M,Ri^!F)$，现在计算 $Ri_*(G)_{\bar x}$。为此，与（第 I 讲 3.2.1）一样，引入 $Y$ 在 $X$ 中的\Term{schematic-closure}{概形闭包} $Z$，得到\Term{immersion}{浸入}的\Term{cartesian-diagram}{笛卡尔图}
\[
\begin{tikzcd}
Y\arrow[r]\arrow[d,"j"']&D(f)\arrow[d]\\
Z\arrow[r,"k"']&X.
\end{tikzcd}
\]
\SourcePage{77}
因而有 $Ri_*(G)\simeq k_*Rj_*(G)$。当然可以假定 $x\in Z$，否则显然有 $Ri_*(G)_{\bar x}=0$。现在引入 $Z$ 在 $\bar x$ 处的\Term{strict-localization}{严格局部化}
$\bar Z=\operatorname{Spec}(\mathcal O_{Z,\bar x})$，并以 $\bar f$ 表示 $f$ 在 $\mathcal O_{Z,\bar x}$ 中的像，记
\[
\bar Y=\bar Z\times_ZY=\bar Z\cap D(\bar f),
\]
以 $\bar G$ 表示 $G$ 在 $\bar Y$ 上的\Term{inverse-image}{逆像}。由（SGA 4 VIII 5.2），有
\[
Ri_*(G)_{\bar x}=Rj_*(G)_{\bar x}\simeq R\Gamma(\bar Y,\bar G).
\]
于是，条件 $\mathrm{(cdloc)}_n$ 通过一个\Term{spectral-sequence}{谱序列}论证给出
\[
R^pi_*(G)_{\bar x}=0
\quad\bigl(p>2d+\operatorname{dim.inj}(F_{\bar x})+\dim(\mathcal O_{Z,x})\bigr).
\]
而且有（参见第 I 讲 3.2.1 的证明）
\[
\dim(\mathcal O_{Z,x})=\dim(\mathcal O_{X,x})-d\geq0,
\]
所以
\[
R^pi_*(G)_{\bar x}=0
\quad\bigl(p>d+\dim(\mathcal O_{X,x})+\operatorname{dim.inj}(F_{\bar x})\bigr),
\]
从而当 $p>2\dim(\mathcal O_{X,x})+\operatorname{dim.inj}(F_{\bar x})$ 时更是如此，这就证明了（*）。

还须证明上界（*）是最优的。置 $f=\operatorname{dim.inj}(F_{\bar x})$，并取一个\Term{finite-type}{有限型} $A_{\bar x}$ 模 $M'$，使
\[
\operatorname{Ext}^f(M',F_{\bar x})\ne0.
\]
设 $Z$ 为 $X$ 中作为 $x$ 的\Term{closure}{闭包}的\Term{integral}{整}闭\Term{subprescheme}{子预概形}，$Y=U\cap Z$ 为 $Z$ 的一个包含 $x$ 的\Term{regular}{正则}\Term{open-subset}{开集}，且它在每一点处的\Term{codimension}{余维}均为 $d=\dim(\mathcal O_{X,x})$。适当缩小 $Y$ 后，存在\Term{locally-constant}{局部常值}\Term{constructible}{可构造} $A_Y$ 模 $M$，使 $M_{\bar x}=M'$。应用\Term{purity-theorem}{纯性定理}，即得
\[
\mathcal Ext^{2d+f}(M,F|_U)_{\bar x}
\simeq\operatorname{Ext}^f(M',F_{\bar x})\ne0,
\]
从而完成证明。
\end{proof}

\SourcePage{78}
注意，若 $F$ 为 $A_X$ 模且 $F_{\bar x}$ 为\Term{injective}{内射}的，则在证明的最后一部分可以取 $M=A_Y$。因此得到：
\begin{SGAresult}{推论}{5.1.1}
在（5.1）的假设下，设 $F$ 为\Term{locally-constant}{局部常值} $A_X$ 模，且 $F_{\bar x}$ 为\Term{injective}{内射}的。则有
\[
\operatorname{dim.q.inj}_{\bar x}(F)
=\operatorname{dimstop}_{\bar x}(F)=2\dim(\mathcal O_{X,x}),
\]
其中使用第 2 节的记号。
\end{SGAresult}
\textbf{注记 5.1.2.}
\textbf{a)} 使用（5.1）的记号，若 $F_{\bar x}=0$，则 $F$ 在 $x$ 的某个\Term{neighborhood}{邻域}内为 $0$，所以对每个 $E\in\operatorname{ob}D^-(X)$ 都有 $R\mathcal Hom(E,F)_{\bar x}=0$。

\textbf{b)} 若 $\operatorname{dim.inj}(F_{\bar x})=+\infty$，则 $\operatorname{dim.q.inj}_{\bar x}(F)=+\infty$；事实上，对每个 $p$，存在\Term{finite-type}{有限型} $A_{\bar x}$ 模 $M'$，使 $\operatorname{Ext}^p(M',F_{\bar x})\ne0$，因而存在一个在 $x$ 的\Term{neighborhood}{邻域}内\Term{locally-constant}{局部常值}的\Term{constructible}{可构造} $A_X$ 模 $M$，使 $\mathcal Ext^p(M,F)_{\bar x}\ne0$。

\SGAsection{6}{\texorpdfstring{$R\mathcal Hom$}{RHom} 的可构造性}
将\Term{ringed-sheaf}{环层} $(\mathbf Z/n\mathbf Z)_X$ 换成一个交换、\Term{locally-constant}{局部常值}且\Term{geometric-fiber}{几何纤维}为\Term{noetherian-ring}{诺特环}的\Term{ringed-sheaf}{环层}时，（第 I 讲 3.3.1）的 a) 中各条件的\Term{equivalence}{等价}性仍然成立：证明完全无须改动。

在具有\Term{resolution-of-singularities}{奇点消解}和\Term{strong-cohomological-purity}{强意义的纯性}（4.4）时，上述条件成立；特别地，当 $X$ 的\Term{dimension}{维数} $\leq1$，或者 $X$ 为零\Term{characteristic}{特征}\Term{excellent}{优良}\Term{prescheme}{预概形}，或者 $X$ 在某个域上\Term{local-scheme}{局部}\Term{finite-type}{有限型}且\Term{dimension}{维数} $\leq2$ 时，这些条件成立。

\SourcePage{79}
\SGAsection{7}{对偶化复形的存在性}
本节中，$X$ 表示配备了一个交换、\Term{locally-constant}{局部常值}且\Term{geometric-fiber}{几何纤维}为\Term{noetherian-ring}{诺特环}的\Term{ringed-sheaf}{环层} $A_X$ 的\Term{prescheme}{预概形}。

\begin{SGAresult}{\Term{lemma}{引理}}{7.1}
设 $F\in\operatorname{ob}D^-_{\mathrm{coh}}(X)$，$G\in\operatorname{ob}D^+_{\mathrm{coh}}(X)$，则
\[
R\mathcal Hom(F,G)\in D^+_{\mathrm{coh}}(X)
\]
（记号见（1.2））。
\end{SGAresult}
\begin{proof}
这是标准的。利用\Term{spectral-sequence}{谱序列}归结到 $F$、$G$ 都集中于\Term{degree}{次数} $0$ 的情形，然后仿照（EGA $0_{\mathrm{III}}$ 12.3.3）进行。
\end{proof}

\begin{SGAresult}{\Term{lemma}{引理}}{7.2}
设 $g:X'\to X$ 为\Term{prescheme}{预概形}的\Term{morphism}{态射}，在 $X'$ 上配备\Term{ringed-sheaf}{环层} $A_{X'}=g^*(A_X)$。则对 $F\in\operatorname{ob}D^-_{\mathrm{coh}}(X)$、$G\in\operatorname{ob}D^+(X)$，\Term{canonical}{典范}\Term{arrow}{箭头}
\[
g^*R\mathcal Hom(F,G)\longrightarrow R\mathcal Hom(g^*F,g^*G)
\]
为\Term{isomorphism}{同构}。
\end{SGAresult}
\begin{proof}
这也是标准的：通过\Term{resolution}{分解}归结到 $F=A_X$（参见 EGA $0_{\mathrm{III}}$ 12.3.5）。
\end{proof}
\textbf{注记 7.2.1.} 前面的\Term{lemma}{引理}是（SGA 6 I）中所载更一般断言的特例。

\begin{SGAresult}{\Term{lemma}{引理}}{7.3}
假定 $X$ 为\Term{quasi-compact}{拟紧}的，$K\in\operatorname{ob}D^+_{\mathrm{coh}}(X)$，且对 $X$ 的每个\Term{geometric-point}{几何点} $\bar x$，$K_{\bar x}$ 均为 [H] V 2.1 意义下的\Term{dualizing-complex}{对偶化复形}。则对每个 $F\in\operatorname{ob}D^b_{\mathrm{coh}}(X)$，有
\[
R\mathcal Hom(F,K)\in\operatorname{ob}D^b_{\mathrm{coh}}(X),
\]
而且 $F$ 关于 $K$ 是\Term{reflexive}{自反}的。
\end{SGAresult}
\SourcePage{80}
\begin{proof}
当然只须对集中于\Term{degree}{次数} $0$ 的 $F$ 验证。由（7.1）有 $R\mathcal Hom(F,K)\in\operatorname{ob}D^+_{\mathrm{coh}}(X)$。来证明 $R\mathcal Hom(F,K)\in D^b_{\mathrm{coh}}(X)$。设 $\bar x$ 为 $X$ 的\Term{geometric-point}{几何点}。由（7.2）有
\[
R\mathcal Hom(F,K)_{\bar x}\simeq R\operatorname{Hom}(F_{\bar x},K_{\bar x}).
\]
若置 $\operatorname{dim.inj}(K_{\bar x})=N(\bar x)$，便有
\[
\mathcal Ext^i(F,K)_{\bar x}=0\qquad(i>N(\bar x)).
\]
由于各个 $\mathcal Ext^i(F,K)$ 都是\Term{locally-constant}{局部常值}的，因而存在 $x$ 的\Term{open-neighborhood}{开邻域} $U$，使当 $i>N(\bar x)$ 时 $\mathcal Ext^i(F,K)|_U=0$；利用 $X$ 的\Term{quasi-compact}{拟紧}性假设即得结论。还须验证\Term{canonical}{典范}\Term{arrow}{箭头}
\[
F\longrightarrow R\mathcal Hom\bigl(R\mathcal Hom(F,K),K\bigr)
\]
为\Term{isomorphism}{同构}。为此只须在 $X$ 的一个\Term{geometric-point}{几何点} $\bar x$ 处考察。在\Term{commutative-diagram}{交换图}
\[
\begin{tikzcd}[column sep=large]
F_{\bar x}\arrow[r]\arrow[d,"\mathrm{id}"']&
R\mathcal Hom\bigl(R\mathcal Hom(F,K),K\bigr)_{\bar x}\arrow[d]\\
F_{\bar x}\arrow[r]&R\operatorname{Hom}\bigl(R\operatorname{Hom}(F_{\bar x},K_{\bar x}),K_{\bar x}\bigr)
\end{tikzcd}
\]
中，右侧竖\Term{arrow}{箭头}由（7.2）为\Term{isomorphism}{同构}，下面的横\Term{arrow}{箭头}由于 $K_{\bar x}$ 为\Term{dualizing-complex}{对偶化复形}也是\Term{isomorphism}{同构}；故上面的横\Term{arrow}{箭头}是\Term{isomorphism}{同构}，证毕。
\end{proof}

\begin{SGAresult}{\Term{lemma}{引理}}{7.4}
假定 $X$ 为\Term{noetherian}{诺特}\Term{regular}{正则}\Term{prescheme}{预概形}，$A_X$ 被某个与 $X$ 的\Term{residue-characteristic}{剩余特征}互素的整数 $n>0$ \Term{annihilate}{零化}，
\SourcePage{81}
并且 $X$ 上的\Term{purity-theorem}{纯性定理}在强意义下成立（4.4）。设 $D$ 为 $X$ 上的\Term{normal-crossings-divisor}{正规交叉除子}（第 I 讲 3.1.5 a)），以 $j:Y=V(D)\to X$ 和 $i:X-Y=U\to X$ 表示\Term{canonical}{典范}包含。设 $K\in\operatorname{ob}D^+_{\mathrm{coh}}(X)$，且对 $X$ 的每个\Term{geometric-point}{几何点} $\bar x$，$K_{\bar x}$ 均为环 $A_{\bar x}$ 的\Term{dualizing-complex}{对偶化复形}（在 [H] V 2.1 的意义下）。则对每个 $E\in\operatorname{ob}D^b_{\mathrm{coh}}(X)$，有
\[
R\mathcal Hom(i_!i^*(E),K)\in\operatorname{ob}D_c^b(X),
\]
且二元组 $(i_!i^*(E),K)$ 为\Term{bidualizing-pair}{双对偶化二元组}。
\end{SGAresult}
\begin{proof}
可以假定 $E$ 集中于\Term{degree}{次数} $0$，即为\Term{locally-constant}{局部常值}\Term{constructible}{可构造} $A_X$ 模。另一方面，由于问题在 $X$ 上对于\Term{etale}{平展}\Term{topology}{拓扑}是\Term{local-scheme}{局部}的，可以假定：
\begin{enumerate}
\item[(a)] $A_X$ 为值为 $A$ 的\Term{constant-sheaf}{常值环层}，$E$ 为常值模，其值仍滥记作 $E$，即 $E=E_X$；
\item[(b)] $D$ 为严格\Term{normal-crossings-divisor}{正规交叉除子}，即 $D=\sum_{i\in I}\operatorname{div}(s_i)$，其中 $s_i\in\Gamma(X,\mathcal O_X)$。
\end{enumerate}
对 $p=\operatorname{card}(I)$ 归纳。当 $p=0$ 时，断言由（7.3）得到。假定 $I=\{1,\ldots,p\}$，置 $Y_i=V(s_i)$，
\[
Y'=V\Bigl(\sum_{1\leq i\leq p-1}\operatorname{div}(s_i)\Bigr),\qquad U'=X-Y',
\]
得到\Term{immersion}{浸入}的\Term{cartesian-diagram}{笛卡尔图}
\[
\begin{tikzcd}
\varnothing\arrow[r]\arrow[d]&U'\cap Y_p\arrow[r,"i'_p"]\arrow[d,"j'_p"]&Y_p\arrow[d,"j_p"]\\
U\arrow[r]&U'\arrow[r,"i'"']&X.
\end{tikzcd}
\]
考察\Term{exact-sequence}{正合列}
\[
0\longrightarrow E_{U,X}\longrightarrow E_{U',X}
\longrightarrow E_{U'\cap Y_p,X}\longrightarrow0.
\]
\SourcePage{82}
由归纳假设，\Term{theorem}{定理}对二元组 $(E_{U',X},K)$ 成立。由于固定 $K$ 时，使 $R\mathcal Hom(F,K)\in\operatorname{ob}D_c^b(X)$ 且 $(F,K)$ 为\Term{bidualizing-pair}{双对偶化二元组}的 $F$ 在 $D_c(X)$ 中组成一个\Term{triangulated-category}{三角子范畴}，所以只须对二元组 $(E_{U'\cap Y_p,X},K)$ 证明\Term{theorem}{定理}。而有
\[
E_{U'\cap Y_p,X}=j_{p*}i'_{p!}(E_{U'\cap Y_p}),
\]
且由（第 I 讲 1.13），这\Term{equivalence}{等价}于在 $Y_p$ 中对二元组
\[
\bigl(i'_{p!}(E_{U'\cap Y_p}),Rj_p^!(K)\bigr)
\]
证明\Term{theorem}{定理}。但由\Term{purity-theorem}{纯性定理}，有（4.3）
\[
Rj_p^!(K)\simeq j_p^*(K)
\mathbin{\mathop\otimes\limits^L}_{\mathbf Z/n\mathbf Z}
(\mu_n)_{Y_p}^{\otimes-1}[-2].
\]
所以 $Rj_p^!(K)\in\operatorname{ob}D^+_{\mathrm{coh}}(Y_p)$，且对 $Y$ 的每个\Term{geometric-point}{几何点} $\bar y$，$(Rj_p^!(K))_{\bar y}$ 都是\Term{dualizing}{对偶化}的；因而由归纳假设得证。
\end{proof}

\begin{SGAresult}{\Term{theorem}{定理}}{7.5}
作如下假设：$X$ 为\Term{noetherian}{诺特}\Term{regular}{正则}\Term{prescheme}{预概形}；$A_X$ 被某个与 $X$ 的\Term{residue-characteristic}{剩余特征}互素的整数 $n>0$ \Term{annihilate}{零化}；$X$ 是\Term{strongly-desingularizable}{强可消解奇点}的（第 I 讲 3.1.5），满足条件 $\mathrm{(reg)}$（第 I 讲 3.1.1）和 $\mathrm{(cdloc)}_n$（第 I 讲 3.1.3）；在 $X$ 上\Term{smooth}{光滑}的每个\Term{prescheme}{预概形}上，\Term{purity-theorem}{纯性定理}都在强意义下成立（4.4）。（若 $X$ 为零\Term{characteristic}{特征}\Term{excellent}{优良}\Term{prescheme}{预概形}，或其\Term{dimension}{维数} $\leq2$ 且在完全域上\Term{smooth}{光滑}，则这些条件满足。）

于是，若 $K\in\operatorname{ob}D^+_{\mathrm{coh}}(X)$，且对 $X$ 的每个\Term{geometric-point}{几何点} $\bar x$，$K_{\bar x}$ 都是\Term{dualizing}{对偶化}的（在 [H] V 2.1 的意义下），则 $K$ 是\Term{dualizing-complex}{对偶化复形}。
\end{SGAresult}
\begin{proof}
须证明 $K$ 满足（第 I 讲 1.7）的条件 (i)--(iii)。对 $X$ 的每个\Term{geometric-point}{几何点} $\bar x$，$K_{\bar x}$ 的\Term{injective-dimension}{内射维数}\Term{finite}{有限}；
\SourcePage{83}
又因为 $K\in\operatorname{ob}D^+_{\mathrm{coh}}(X)$，函数 $x\mapsto\operatorname{dim.inj}(K_{\bar x})$ 为\Term{locally-constant}{局部常值}的，由 $X$ 为\Term{noetherian}{诺特}的知其\Term{bounded}{有界}。因而由（5.1）得 $K$ 的\Term{quasi-injective-dimension}{拟内射维数}\Term{finite}{有限}。

由第 6 节的讨论，上引之处的条件 (ii) 也满足。

还须证明条件 (iii)，即每个 $F\in\operatorname{ob}D_c^b(X)$ 关于 $K$ 都是\Term{reflexive}{自反}的。通过\Term{resolution}{分解}，可以假定 $F$ 集中于\Term{degree}{次数} $0$。另一方面，由于问题在 $X$ 上对于\Term{etale}{平展}\Term{topology}{拓扑}是\Term{local-scheme}{局部}的，可以假定 $A_X$ 为值为 $A$ 的\Term{constant-sheaf}{常值环层}。于是\Term{resolution}{分解}（1.3）使我们归结到 $F=f_*i_!(E_U)$，其中 $f:Y\to X$ 为\Term{finite-morphism}{有限态射}，$Y$ 为\Term{integral}{整}的，其\Term{generic-point}{一般点}在其像上\Term{separable}{可分}；$i:U\to Y$ 为\Term{nonempty}{非空}\Term{regular}{正则}\Term{open-subset}{开集}，$E_U$ 为常值 $A_U$ 模，其值 $E$ 是\Term{finite-type}{有限型} $A$ 模。由于 $X$ 为\Term{strongly-desingularizable}{强可消解奇点}的，存在（第 I 讲 3.1.5）\Term{projective-morphism}{射影}\Term{birational-morphism}{双有理}\Term{morphism}{态射} $h:Y'\to Y$，使 $Y'$ 为\Term{regular}{正则}的，$h$ \Term{induction-isomorphism}{诱导同构} $h^{-1}(U)\to U$，且 $h^{-1}(U)$ 为 $Y'$ 中某个\Term{normal-crossings-divisor}{正规交叉除子}的补集。借助 $h$ 将 $h^{-1}(U)$ 与 $U$ 等同，就得到\Term{commutative-diagram}{交换图}
\[
\begin{tikzcd}[column sep=large]
&Y'\arrow[d,"h"]\arrow[dr,"g"]&\\
U\arrow[ur,"i'"]\arrow[r,"i"']&Y\arrow[r,"f"']&X.
\end{tikzcd}
\]
从而有
\begin{align*}
f_*i_!(E_U)&\simeq f_*Rh_*i'_!(E_U)\\
&\simeq Rg_*i'_!(E_U)\\
&\simeq Rg_*i'_!i'^*(E_{Y'}).
\end{align*}
由（第 I 讲 1.13），只须对二元组
\SourcePage{84}
\[
\bigl(i'_!i'^*(E_{Y'}),R^!g(K)\bigr)
\]
证明\Term{biduality}{双对偶性}。但将 $Y'$ \Term{local-scheme}{局部}地\Term{embedding}{嵌入}到一个在 $X$ 上\Term{smooth}{光滑}的\Term{prescheme}{预概形} $X'$ 中，再在 $X'$ 上应用\Term{strong-cohomological-purity}{强意义的纯性}（4.4），就得到 $R^!g(K)$ \Term{local-scheme}{局部}地\Term{isomorphism}{同构}于 $g^*(K)$，差一个\Term{degree}{次数}\Term{translation}{平移}。因此 $R^!g(K)\in\operatorname{ob}D^+_{\mathrm{coh}}(Y')$，且对 $Y'$ 的每个\Term{geometric-point}{几何点} $\bar y$，$R^!g(K)_{\bar y}$ 都是\Term{dualizing}{对偶化}的。此时只须应用（7.4）。

此外，前面的证明也证明了
\[
R\mathcal Hom(F,K)\in\operatorname{ob}D_c^b(X)
\]
（这里使用\Term{proper}{固有}\Term{morphism}{态射}的\Term{finiteness-theorem}{有限性定理}（SGA 4 XIV））；这可以使不信任第 6 节那番议论的读者放心。

证明至此完成。
\end{proof}

\bigskip
\begin{center}\textbf{参考文献}\end{center}
\noindent [1]\quad Verdier, J.-L.，《\Term{locally-compact-space}{局部紧空间}的\Term{dimension}{维数}》（\textit{Dimension des espaces localement compacts}），\textit{C. R. Acad. Sci. Paris}，261，第 5293--5296 页（1965）。（另见 Verdier, J.-L.，《\Term{locally-compact-space}{局部紧空间}\Term{cohomology}{上同调}中的\Term{duality}{对偶性}》（\textit{Dualité dans la cohomologie des espaces localement compacts}），\textit{Séminaire Bourbaki}，第 300 讲，1965--66。）

\medskip
\noindent SGA A = SGA 4。
